\documentclass[11pt,a4paper]{article}
\usepackage[T1]{fontenc}
\usepackage[utf8]{inputenc}
\usepackage[margin=23mm,headheight=15pt]{geometry}
\usepackage{amsmath,amssymb,mathtools,amsthm}
\usepackage{newpxtext,newpxmath}
\usepackage{microtype,xcolor,enumitem,environ,needspace,fancyhdr,etoolbox}
\usepackage[colorlinks=true,linkcolor=heading,urlcolor=heading]{hyperref}
\definecolor{heading}{HTML}{244B65}
\definecolor{hintcolor}{HTML}{765638}
\setlength{\parindent}{0pt}
\setlength{\parskip}{5pt plus 1pt}
\setlength{\emergencystretch}{2em}
\setlist{topsep=4pt,itemsep=4pt,parsep=2pt,leftmargin=*}
\allowdisplaybreaks[1]
\newtheoremstyle{exostyle}{8pt}{6pt}{\normalfont}{}{\color{heading}\bfseries}{.}{\newline}{}
\theoremstyle{exostyle}
\newtheorem{exercise}{Exercise}
\BeforeBeginEnvironment{exercise}{\par\addvspace{10pt}\Needspace{9\baselineskip}%
  \begingroup\color{heading!45}\hrule height .6pt\endgroup\nobreak}
\newif\ifcorr
\ifdefined\StatementOnly\corrfalse\else\corrtrue\fi
\newcommand{\showsolutions}{\corrtrue}
\newcommand{\hidesolutions}{\corrfalse}
\NewEnviron{hints}{\ifcorr\par\Needspace{4\baselineskip}\noindent
  {\color{hintcolor}\bfseries Useful hints}\par\BODY\par\medskip\fi}
\NewEnviron{solution}{\ifcorr\par\Needspace{4\baselineskip}\noindent
  {\color{heading}\bfseries Detailed solution}\par\BODY\par\medskip\fi}
\newcommand{\R}{\mathbb{R}}
\newcommand{\push}{\#}
\newcommand{\W}{\mathcal{W}}
\newcommand{\Wone}{\mathcal{W}_1}
\newcommand{\Wtwo}{\mathcal{W}_2}
\newcommand{\Unif}{\mathcal U}
\newcommand{\ip}[2]{\langle #1,#2\rangle}
\newcommand{\gradW}{\mathord{\nabla\mkern-9.3mu\nabla}}
\pagestyle{fancy}
\fancyhf{}
\fancyhead[L]{\small\color{heading}Optimal Transport / Exercises}
\fancyhead[R]{\small\ifcorr Solutions\else Statements\fi}
\fancyfoot[C]{\small\thepage}
\renewcommand{\headrulewidth}{0.3pt}
\renewcommand{\maketitle}{\thispagestyle{plain}\begingroup
  \color{heading}\LARGE\bfseries\SheetTitle\par\endgroup
  \vspace{3pt}{\large\ifcorr Statements, hints and worked solutions\else Exercise statements\fi}\par
  \vspace{7pt}\hrule\vspace{9pt}
  \ifcorr Each exercise is followed by question-by-question hints and a worked solution.
  Try the questions before reading the hints.\else
  This version contains the exercises only. Hints and worked solutions are available in the companion PDF.\fi
  \par\medskip}
\hypersetup{pdfauthor={Optimal Transport for Machine Learners},pdfsubject={Exercises in optimal transport}}

\newcommand{\SheetTitle}{Sliced Optimal Transport}
\DeclareMathOperator{\SW}{\mathcal{SW}}
\DeclareMathOperator{\Law}{Law}
\newcommand{\EE}{\mathbb E}
\newcommand{\PPtwo}{\mathcal P_2}
\newcommand{\Couplings}{\mathcal U}
\newcommand{\Sphere}{\mathbb S}
\newcommand{\Gauss}{\mathcal N}
\hypersetup{pdftitle={Sliced Optimal Transport -- Exercises},
  pdfauthor={Kimia Nadjahi (original exercises); OT4ML (adaptation and solutions)}}

\begin{document}
\maketitle
\textbf{Original exercises: Kimia Nadjahi.} Adapted for OT4ML.

\textbf{Notation.}
Let $\PPtwo(\R^d)$ be the probability measures with a finite second moment
and $\Couplings(\alpha,\beta)$ the couplings of $\alpha$ and $\beta$.
Thus $\Wtwo(\alpha,\beta)^2$ is the infimum of
$\int |x-y|^2\,d\pi(x,y)$ over $\pi\in\Couplings(\alpha,\beta)$.
For $\theta\in\Sphere^{d-1}$, write
$P_\theta(x)=\ip{\theta}{x}$ and $\alpha_\theta=(P_\theta)_\push\alpha$.
The order-two sliced Wasserstein distance is defined by
\[
  \SW_2(\alpha,\beta)^2
    =\int_{\Sphere^{d-1}}\Wtwo(\alpha_\theta,\beta_\theta)^2
      \,d\varsigma(\theta),
\]
where $\varsigma$ is the uniform \emph{probability} measure on the sphere.

\begin{exercise}[Comparing centered isotropic Gaussians]\label{ex-sliced-gaussians}
Let $\alpha=\Gauss(0,\sigma^2 I_d)$ and $\beta=\Gauss(0,\tau^2 I_d)$ with
$\sigma,\tau\geq0$, where a zero standard deviation denotes $\delta_0$.
\begin{enumerate}[label=(\alph*)]
\item Compute $\SW_2(\alpha,\beta)$.
\item Compute $\Wtwo(\alpha,\beta)$ and compare the two distances.
\end{enumerate}
\end{exercise}
\begin{hints}
\begin{enumerate}[label=(\alph*)]
\item If $G\sim\Gauss(0,I_d)$, what is the variance of $\ip{\theta}{G}$?
  In one dimension, try coupling the two Gaussians with the same standard
  Gaussian. To prove optimality, expand the squared cost and bound the
  cross moment by Cauchy--Schwarz.
\item In $\R^d$, use the same standard Gaussian \emph{vector} for the
  two marginals. For any other coupling, apply Cauchy--Schwarz to its
  expected inner product.
\end{enumerate}
\end{hints}
\begin{solution}
\textbf{(a) Project, transport on the line, then average.}
Write $G=(G_1,\ldots,G_d)\sim\Gauss(0,I_d)$.
The projection $\ip{\theta}{G}=\sum_i\theta_iG_i$ is Gaussian. Its mean
is zero and its variance is
$\sum_{i,j}\theta_i\theta_j\EE[G_iG_j]=\sum_i\theta_i^2=1$.
Since $\sigma G$ and $\tau G$ have laws $\alpha$ and $\beta$, respectively,
\[
  \alpha_\theta=\Gauss(0,\sigma^2),\qquad
  \beta_\theta=\Gauss(0,\tau^2).
\]
For any coupling $(S,T)$ of these projected laws,
\[
  \EE[(S-T)^2]=\sigma^2+\tau^2-2\EE[ST]
        \geq\sigma^2+\tau^2-2\sigma\tau=(\sigma-\tau)^2.
\]
The inequality uses $\EE[ST]\leq(\EE S^2)^{1/2}(\EE T^2)^{1/2}$.
It is attained by $(S,T)=(\sigma Z,\tau Z)$ with
$Z\sim\Gauss(0,1)$, since $\EE[ST]=\sigma\tau\EE Z^2=\sigma\tau$.
This proof needs no division by $\sigma$ or $\tau$, so it also covers
a Dirac marginal. Every slice has the same squared cost. Integrating
against $\varsigma$, which has total mass one, gives
\[
  \boxed{\SW_2(\alpha,\beta)=|\sigma-\tau|.}
\]

\par\Needspace{10\baselineskip}
\textbf{(b) Transport in the ambient space.}
For a coupling $(X,Y)$ of $\alpha$ and $\beta$, summing the coordinate
variances gives $\EE|X|^2=d\sigma^2$ and $\EE|Y|^2=d\tau^2$. Expanding
the squared norm and applying Cauchy--Schwarz in the space of
square-integrable random vectors yields
\begin{align*}
  \EE|X-Y|^2
    &=d\sigma^2+d\tau^2-2\EE\ip{X}{Y}\\
    &\geq d\sigma^2+d\tau^2
            -2(\EE|X|^2)^{1/2}(\EE|Y|^2)^{1/2}
     =d(\sigma-\tau)^2.
\end{align*}
The coupling $(X,Y)=(\sigma G,\tau G)$ attains the bound because its
displacement is $(\sigma-\tau)G$ and $\EE|G|^2=d$. Thus
\[
  \boxed{\Wtwo(\alpha,\beta)=\sqrt d\,|\sigma-\tau|,\qquad
    \SW_2(\alpha,\beta)=d^{-1/2}\Wtwo(\alpha,\beta).}
\]
The coupling that attains the ambient bound projects to the
one-dimensional optimizer in every direction. Exercise~\ref{ex-sliced-plans}
shows that this property can fail for other measures.
\end{solution}

\begin{exercise}[Projected optimal transport plans]\label{ex-sliced-plans}
Let $\alpha,\beta\in\PPtwo(\R^d)$ and let
$\pi\in\Couplings(\alpha,\beta)$ be optimal for the squared Euclidean
cost. For a unit direction $\theta$, define
$\pi_\theta=(P_\theta,P_\theta)_\push\pi$, where the map sends
$(x,y)$ to $(P_\theta(x),P_\theta(y))$.
\begin{enumerate}[label=(\alph*)]
\item Show that $\pi_\theta\in\Couplings(\alpha_\theta,\beta_\theta)$.
\item Must $\pi_\theta$ be optimal for the squared-distance cost on
  the line? Justify your answer, giving an explicit example if not.
\end{enumerate}
\end{exercise}
\begin{hints}
\begin{enumerate}[label=(\alph*)]
\item If $(X,Y)\sim\pi$, identify the separate laws of
  $(P_\theta(X),P_\theta(Y))$.
\item In $\R^2$, try equal weights on each of $\{u,-u\}$ and
  $\{v,-v\}$. Choose $u,v$ so pairing the vectors with the same signs
  is cheaper in the plane but their first coordinates have opposite signs.
\end{enumerate}
\end{hints}
\begin{solution}
\textbf{(a) Feasibility survives projection.}
If $(X,Y)\sim\pi$, then $\pi_\theta$ is the joint law of
$(P_\theta(X),P_\theta(Y))$. Its first marginal is
$\Law(P_\theta(X))=(P_\theta)_\push\alpha$; its second is
$\Law(P_\theta(Y))=(P_\theta)_\push\beta$. Equivalently, for every Borel
set $A\subset\R$,
\begin{align*}
  \pi_\theta(A\times\R)
    &=\pi(P_\theta^{-1}(A)\times\R^d)=\alpha_\theta(A),\\
  \pi_\theta(\R\times A)
    &=\pi(\R^d\times P_\theta^{-1}(A))=\beta_\theta(A).
\end{align*}
This does not use the optimality of $\pi$.

\textbf{(b) A strict counterexample.}
In $\R^2$ let $u=(1,1)$ and $v=(-1,2)$, and take
$\alpha=\tfrac12(\delta_u+\delta_{-u})$ and
$\beta=\tfrac12(\delta_v+\delta_{-v})$.
With row order $(u,-u)$ and column order $(v,-v)$ the matrix of squared
distances is
\[
  \begin{pmatrix}
    |u-v|^2 & |u+v|^2\\
    |u+v|^2 & |u-v|^2
  \end{pmatrix}
  =\begin{pmatrix}5&9\\9&5\end{pmatrix},
\]
since $u-v=(2,-1)$ and $u+v=(0,3)$.
The marginal constraints force every feasible matrix to be
$\left(\begin{smallmatrix}a&1/2-a\\1/2-a&a\end{smallmatrix}\right)$
for $0\leq a\leq1/2$: once the top-left entry is chosen, its row and
column sums determine the two off-diagonal entries and then the last
entry. Its cost is $2a\cdot5+2(1/2-a)\cdot9=9-8a$.
Hence the unique optimum is $a=1/2$, with cost $5$ and plan
$\pi=\tfrac12\delta_{(u,v)}+\tfrac12\delta_{(-u,-v)}$.

Project onto the first coordinate, $\theta=e_1$. Both projected
marginals are $\tfrac12(\delta_{-1}+\delta_1)$, because
$P_\theta(u)=1$ and $P_\theta(v)=-1$. Nevertheless, the projected plan is
\[
  \pi_\theta=\tfrac12\delta_{(1,-1)}+\tfrac12\delta_{(-1,1)}.
\]
Its cost is $\tfrac12(1-(-1))^2+\tfrac12(-1-1)^2=4$, while the identity
coupling of the two equal projected measures costs zero.
Thus feasibility survives, but optimality can fail even for the
projection of a \emph{unique} ambient optimizer.

For $d>2$ the same example embeds by adding zero coordinates. For $d=1$,
in contrast, the only projections are $x\mapsto x$ and $x\mapsto -x$.
Each is a bijective isometry, so it preserves optimality.

\textbf{Consequence: why slicing gives a lower bound.}
Even when $\pi_\theta$ is not optimal, its feasibility gives
$\Wtwo(\alpha_\theta,\beta_\theta)^2
\leq\int|\ip{\theta}{x-y}|^2\,d\pi(x,y)$.
For any $z\in\R^d$, the spherical average of
$|\ip{\theta}{z}|^2$ is $|z|^2/d$: one can rotate $z$ to $|z|e_1$
and use $\int\theta_1^2\,d\varsigma=1/d$ by symmetry and
$\sum_i\theta_i^2=1$. Averaging the inequality and swapping integrals
(the integrand is nonnegative) proves
\[
  \SW_2(\alpha,\beta)^2
     \leq \frac1d\int|x-y|^2\,d\pi(x,y)
     =\frac1d\Wtwo(\alpha,\beta)^2.
\]
For the example above this inequality is strict, not merely a failure
in an isolated direction. For $\theta$ near $e_1$, the projected
values $a_\theta=P_\theta(u)$ and $b_\theta=P_\theta(v)$ still have
opposite signs; switching the projected pairing reduces the cost
from $(a_\theta-b_\theta)^2$ to $(a_\theta+b_\theta)^2$.
This open set of directions has positive $\varsigma$-measure.
\end{solution}

\begin{exercise}[Translation decomposition for $\SW_2$]\label{ex-sliced-centering}
Let $\alpha,\beta\in\PPtwo(\R^d)$. Write $m_\alpha=\int x\,d\alpha(x)$
and $m_\beta=\int y\,d\beta(y)$ for their means, and
$\alpha_0,\beta_0$ for their centered translates.
\begin{enumerate}[label=(\alph*)]
\item Compute the means of $\alpha_\theta$ and $\beta_\theta$ for
  each fixed $\theta\in\Sphere^{d-1}$.
\item Find the unique translations that center $\alpha$ and $\beta$.
  Show that centering and projecting commute:
  $(\alpha_\theta)_0=(P_\theta)_\push\alpha_0$ and
  $(\beta_\theta)_0=(P_\theta)_\push\beta_0$.
\item Prove that
\[
  \SW_2(\alpha,\beta)^2
    =\frac1d|m_\alpha-m_\beta|^2+\SW_2(\alpha_0,\beta_0)^2.
\]
\end{enumerate}
\end{exercise}
\begin{hints}
\begin{enumerate}[label=(\alph*)]
\item Use linearity of $P_\theta$ and the definition of the pushforward.
\item Translating a random vector by $-m$ subtracts $m$ from its mean.
  Compute $P_\theta(x-m)$.
\item On the line, expand the squared difference after separately
  centering the two coordinates of a coupling. Which term drops out
  for \emph{every} coupling? Finally use
  $\int\theta\theta^\top\,d\varsigma(\theta)=I_d/d$.
\end{enumerate}
\end{hints}
\begin{solution}
\textbf{(a) A projected mean is the projection of the mean.}
Finite second moments imply finite first moments by Cauchy--Schwarz.
For a fixed direction, the pushforward formula gives
\[
  \int s\,d\alpha_\theta(s)
     =\int P_\theta(x)\,d\alpha(x)
     =\ip{\theta}{\int x\,d\alpha(x)}
     =\ip{\theta}{m_\alpha}.
\]
Replacing $\alpha$ by $\beta$ gives $\int t\,d\beta_\theta(t)=
\ip{\theta}{m_\beta}$. There is no average over directions at this step.

\textbf{(b) Center first or project first.}
For any $u\in\R^d$, the mean of $(x\mapsto x-u)_\push\alpha$ is
$\int(x-u)\,d\alpha(x)=m_\alpha-u$, using the unit mass of $\alpha$.
It is zero precisely when $u=m_\alpha$. Likewise, $\beta$ is centered
by subtracting $m_\beta$, and no other vector. If $X\sim\alpha$, then
\[
  P_\theta(X-m_\alpha)
    =P_\theta(X)-\ip{\theta}{m_\alpha}.
\]
The left-hand side has law $(P_\theta)_\push\alpha_0$ and the
right-hand side has law $(\alpha_\theta)_0$ by part~(a).
The same calculation with $\beta$ gives the second identity.

\par\Needspace{11\baselineskip}
\textbf{(c) Split the cost in each direction.}
First let $\eta,\zeta\in\PPtwo(\R)$ have means $a,b$, and let
$(S,T)$ be \emph{any} coupling of them. Write $S_0=S-a$ and $T_0=T-b$.
Then $S-T=(S_0-T_0)+(a-b)$ and
\begin{align*}
  \EE[(S-T)^2]
    &=\EE[(S_0-T_0)^2]+2(a-b)\EE[S_0-T_0]+(a-b)^2\\
    &=\EE[(S_0-T_0)^2]+(a-b)^2,
\end{align*}
since $\EE S_0=\EE T_0=0$ regardless of their dependence.
Subtracting $a,b$ from the two coordinates is a bijection between
the couplings of $\eta,\zeta$ and those of $\eta_0,\zeta_0$:
the inverse simply adds $a,b$ back. The constant term does not depend
on the chosen coupling, so taking the infimum gives
$\Wtwo(\eta,\zeta)^2=(a-b)^2+\Wtwo(\eta_0,\zeta_0)^2$.
Apply this to $\alpha_\theta,\beta_\theta$ and use parts (a) and (b):
\[
  \Wtwo(\alpha_\theta,\beta_\theta)^2
    =|\ip{\theta}{m_\alpha-m_\beta}|^2+
      \Wtwo\bigl((P_\theta)_\push\alpha_0,(P_\theta)_\push\beta_0\bigr)^2.
\]

\textbf{Average the squared mean displacement.}
For completeness, the matrix $M=\int\theta\theta^\top\,d\varsigma(\theta)$
has vanishing off-diagonal entries: flipping the sign of one
coordinate preserves $\varsigma$ and negates the corresponding
product of distinct coordinates. Its diagonal entries are equal by
permutation symmetry and add to $\int|\theta|^2\,d\varsigma=1$.
Thus $M=I_d/d$, including $d=1$. Setting $w=m_\alpha-m_\beta$ gives
\[
  \int|\ip{\theta}{w}|^2\,d\varsigma(\theta)
     =w^\top M w=\frac{|w|^2}{d}.
\]
Averaging the directional identity now yields
\[
  \boxed{\SW_2(\alpha,\beta)^2
    =\frac{|m_\alpha-m_\beta|^2}{d}+\SW_2(\alpha_0,\beta_0)^2.}
\]
For example, if $\beta=(x\mapsto x+h)_\push\alpha$, then
$m_\beta=m_\alpha+h$ and $\beta_0=\alpha_0$, so
$\SW_2(\alpha,\beta)=|h|/\sqrt d$. If the means already agree, the first
term vanishes and the distance is entirely between the centered laws.
\end{solution}

\par\addvspace{7pt}\Needspace{3\baselineskip}
\textbf{Further reading.}
See the sections on sliced Wasserstein (including Gaussians) and on
the Euclidean mean decomposition in
\href{https://www.gpeyre.com/ot4ml/OT4ML/OT4ML.pdf}
     {\emph{Optimal Transport for Machine Learners}}.
\end{document}
